% GENERATED FILE. Edit canonical lessons, then run npm run generate:books.
% canonical-source-hash: fnv1a64:ff04929106c65879
% canonical-lessons: FR-C437-impasse, FR-C437-ruelle, FR-C437-avenue, FR-C437-boulevard, FR-C437-passant

\chapter{Petites rues et grandes avenues}
\label{ch:fr-petites-rues-et-grandes-avenues}
\hlchaptermodality{437}

\begin{chapteropening}
\textbf{By the end of this chapter:} \emph{I can describe streets: the dead end, the alley, the avenue, the boulevard and the passer-by.}

Complete the last lesson of chapter 437: I can describe streets: the dead end, the alley, the avenue, the boulevard and the passer-by.
\end{chapteropening}

\section[l'impasse]{l'impasse — the dead end; the deadlock}

\label{lesson:FR-C437-impasse}



\begin{quote}
Before the new one: say the French for to go round, then the French for to turn back.
\end{quote}

\subsection*{You'll want to know: l'impasse}
\textbf{l'impasse} — \textquotedblleft{}the dead end; the deadlock\textquotedblright{}.

\begin{quote}
L'hôtel est au fond d'une impasse calme.
\end{quote}

\begin{quote}
\emph{The hotel is at the end of a quiet dead end.}
\end{quote}

A feminine noun. Figurative: \emph{être dans l'impasse} — to be at a deadlock.

\begin{cousinweb}
\emph{impasse} is built from \emph{in-} 'not' and \emph{passe}, from \emph{passer} 'to pass' (EtymDB), which goes back to Vulgar Latin \emph{passāre}, from \emph{passus} 'a step'.
\end{cousinweb}

\begin{grammarlens}[title={What you've built}]
One more word for telling, explaining and describing at B1.
\end{grammarlens}

\subsection*{Your turn}
\begin{itemize}
\raggedright
  \item \emph{Say it:} \emph{l'impasse}
  \item \emph{Say it:} \emph{l'impasse}, once more
  \item \emph{Say it:} a sentence of your own with \emph{l'impasse}
  \item \emph{Recall:} say the French for to go round, then the French for to turn back, then say \emph{l'impasse} again
\end{itemize}

\begin{tcolorbox}[breakable,colback=teal!4,colframe=teal!35!black,title={Before you move on}]
What does l'impasse mean? (\textquotedblleft{}The dead end; the deadlock\textquotedblright{}.) Say it once more.
\end{tcolorbox}
\section[la ruelle]{la ruelle — the alley, the lane}

\label{lesson:FR-C437-ruelle}



\begin{quote}
Before the new one: say the French for to turn back, then the French for the dead end.
\end{quote}

\subsection*{You'll want to know: la ruelle}
\textbf{la ruelle} — \textquotedblleft{}the alley, the lane\textquotedblright{}.

\begin{quote}
Nous nous promenons dans les ruelles du vieux quartier.
\end{quote}

\begin{quote}
\emph{We stroll through the alleys of the old quarter.}
\end{quote}

A small \emph{rue}.

\begin{cousinweb}
\emph{ruelle} is a small \emph{rue} 'street', which is inherited from Latin \emph{rūga} 'wrinkle, fold' (Lewis \& Short). EtymDB lists it without a source.
\end{cousinweb}

\begin{grammarlens}[title={What you've built}]
One more word for telling, explaining and describing at B1.
\end{grammarlens}

\subsection*{Your turn}
\begin{itemize}
\raggedright
  \item \emph{Say it:} \emph{la ruelle}
  \item \emph{Say it:} \emph{la ruelle}, once more
  \item \emph{Say it:} a sentence of your own with \emph{la ruelle}
  \item \emph{Recall:} say the French for to turn back, then the French for the dead end, then say \emph{la ruelle} again
\end{itemize}

\begin{tcolorbox}[breakable,colback=teal!4,colframe=teal!35!black,title={Before you move on}]
What does la ruelle mean? (\textquotedblleft{}The alley, the lane\textquotedblright{}.) Say it once more.
\end{tcolorbox}
\section[l'avenue]{l'avenue — the avenue}

\label{lesson:FR-C437-avenue}



\begin{quote}
Before the new one: say the French for the dead end, then the French for the alley.
\end{quote}

\subsection*{You'll want to know: l'avenue}
\textbf{l'avenue} — \textquotedblleft{}the avenue\textquotedblright{}.

\begin{quote}
Les grands magasins sont sur l'avenue principale.
\end{quote}

\begin{quote}
\emph{The big shops are on the main avenue.}
\end{quote}

A feminine noun.

\begin{cousinweb}
\emph{avenue} is the old feminine past participle of \emph{avenir} 'to come to' (EtymDB), from Latin \emph{advenīre}, built on \emph{venīre} 'to come': the way that leads up to a place.
\end{cousinweb}

\begin{grammarlens}[title={What you've built}]
One more word for telling, explaining and describing at B1.
\end{grammarlens}

\subsection*{Your turn}
\begin{itemize}
\raggedright
  \item \emph{Say it:} \emph{l'avenue}
  \item \emph{Say it:} \emph{l'avenue}, once more
  \item \emph{Say it:} a sentence of your own with \emph{l'avenue}
  \item \emph{Recall:} say the French for the dead end, then the French for the alley, then say \emph{l'avenue} again
\end{itemize}

\begin{tcolorbox}[breakable,colback=teal!4,colframe=teal!35!black,title={Before you move on}]
What does l'avenue mean? (\textquotedblleft{}The avenue\textquotedblright{}.) Say it once more.
\end{tcolorbox}
\section[le boulevard]{le boulevard — the boulevard}

\label{lesson:FR-C437-boulevard}



\begin{quote}
Before the new one: say the French for the alley, then the French for the avenue.
\end{quote}

\subsection*{You'll want to know: le boulevard}
\textbf{le boulevard} — \textquotedblleft{}the boulevard\textquotedblright{}.

\begin{quote}
Le bus passe par le boulevard qui fait le tour du centre.
\end{quote}

\begin{quote}
\emph{The bus goes along the boulevard that runs round the centre.}
\end{quote}

Often lined with trees.

\begin{cousinweb}
\emph{boulevard} was borrowed from Middle Dutch \emph{bolwerc} 'bulwark' (EtymDB): first the top of a town wall, later the wide road built where the wall had been.
\end{cousinweb}

\begin{grammarlens}[title={What you've built}]
One more word for telling, explaining and describing at B1.
\end{grammarlens}

\subsection*{Your turn}
\begin{itemize}
\raggedright
  \item \emph{Say it:} \emph{le boulevard}
  \item \emph{Say it:} \emph{le boulevard}, once more
  \item \emph{Say it:} a sentence of your own with \emph{le boulevard}
  \item \emph{Recall:} say the French for the alley, then the French for the avenue, then say \emph{le boulevard} again
\end{itemize}

\begin{tcolorbox}[breakable,colback=teal!4,colframe=teal!35!black,title={Before you move on}]
What does le boulevard mean? (\textquotedblleft{}The boulevard\textquotedblright{}.) Say it once more.
\end{tcolorbox}
\section[le passant]{le passant — the passer-by}

\label{lesson:FR-C437-passant}



\begin{quote}
Before the new one: say the French for the avenue, then the French for the boulevard.
\end{quote}

\subsection*{You'll want to know: le passant}
\textbf{le passant} — \textquotedblleft{}the passer-by\textquotedblright{}.

\begin{quote}
J'ai demandé mon chemin à un passant.
\end{quote}

\begin{quote}
\emph{I asked a passer-by for directions.}
\end{quote}

Feminine \emph{la passante}.

\begin{cousinweb}
\emph{passant} is the present participle of \emph{passer} 'to pass', used as a noun (EtymDB); \emph{passer} goes back to Vulgar Latin \emph{passāre}, from \emph{passus} 'a step'.
\end{cousinweb}

\begin{grammarlens}[title={What you've built}]
One more word for telling, explaining and describing at B1.
\end{grammarlens}

\subsection*{Your turn}
\begin{itemize}
\raggedright
  \item \emph{Say it:} \emph{le passant}
  \item \emph{Say it:} \emph{le passant}, once more
  \item \emph{Say it:} a sentence of your own with \emph{le passant}
  \item \emph{Recall:} say the French for the avenue, then the French for the boulevard, then say \emph{le passant} again
\end{itemize}

\begin{tcolorbox}[breakable,colback=teal!4,colframe=teal!35!black,title={Before you move on}]
What does le passant mean? (\textquotedblleft{}The passer-by\textquotedblright{}.) Say it once more.
\end{tcolorbox}
